\section{Further research questions and proposed directions}
\label{ado:sec:research}
The proofs above suggest a more structured programme than an unrestricted
search for small numerical relations. The following questions separate
analytic geometry, constructive certificates, and arithmetic reductions.

\subsection{Fractional leading indices}
Determine the largest real-index domain on which a finite signed moment
density exists and has exactly $d-1$ sign changes. The first focused
case is depth three with a real first index and integer second index.
The all-positive depth-two results show that absence of a finite signed
measure need not imply failure of zero-freeness. Consequently there are
two different problems: extending the present density theorem, and
extending the zero theorem by a compensated or distributional method.

For real leading indices at least one, a plausible testable target is
that the same $d-1$ oscillation statement persists. This is a proposed
conjectural direction, not a claim supported here by an exhaustive
fractional-parameter experiment. A proof would need a substitute for
the derivative-zero allocation used by $\Aop$.

\subsection{Normalized angular motion and its higher-depth analogue}
Theorem~\ref{ado:thm:radialmotion} settles strict decrease of the angles,
but not the manuscript's finer quotient
$\eta(\rho)=\cos\theta(\rho)/\rho$. At depth two, retain the
existing integer-outer-index normalized-radius conjecture exactly as
stated in the manuscript. At higher depth, a more natural displacement
is $(\alpha_k-\theta_k(\rho))/\rho$. Determine its monotonicity
regions, transition loci, and whether compensator-root interlacing
constrains them. The explicit coefficient in
\eqref{ado:eq:radialexpansion} provides an exact local diagnostic,
not a global conclusion.

\subsection{Motion under changes of an index or the terminal seed}
Raising the first index gives strict interlacing of density zeros.
Does it also order every angular zero in the expected direction toward
$\alpha_k$? Formula \eqref{ado:eq:largeouter} proves the limiting
behavior but does not compare every pair of finite consecutive outer
indices. A positive answer would require comparing both the compensator
and its positive transform, not only the locations of the compensator
roots. Analogous questions for an inner index, the terminal shift
$\lambda$, and positive terminal mixtures are also well posed.

\subsection{Sharper variation norms and faster certificates}
The bound $\norm\kappa_1\le2^{d-1}$ is uniform and easy to certify,
but it is not asserted optimal. Determine the best constant at fixed
depth, or at fixed depth and first index. The exact logistic kernels
supply a useful family in which the variation can be expressed through
one-dimensional integrals between known roots. A smaller certified
variation bound would directly tighten \eqref{ado:eq:centerederror}.
Another target is a full-slit evaluator with comparable exact
certificates and a conformal or Chebyshev expansion adapted to the cut.

\subsection{Constructive compensator roots and positive approximation}
Develop rational interval enclosures for all $r_j$ without assuming
closed forms. The proved interlacing laws can supply nested brackets
when an index or the depth changes. Once such enclosures are available,
combine the positive moment matrices of
\eqref{ado:eq:compensatedmoments} with quadrature or Pad\'e-type
approximants to produce two-sided bounds on real arguments. The
separation between an exact root theorem and a finite root-isolation
certificate should be retained throughout.

\subsection{Large-depth geometry}
For the all-ones family, \eqref{ado:eq:logisticroots} and
\eqref{ado:eq:edgeasymptotic} determine the density-edge scale.
For general index patterns, determine whether the edge roots still
have exponential scale in depth, and how the leading exponents depend
on the index profile. Compare this with the angular boundary layers
near $k=1$ and $k=d-1$. Uniform asymptotics with both $d$ and the
outer index growing would require estimates beyond the fixed-depth
remainder in Theorem~\ref{ado:thm:largeouter}.

\subsection{Radii beyond one and zero-count transitions}
The first target beyond the proved radius range is the explicit all-ones
family. There the problem reduces to the argument of
$-\log(1-\rho e^{\ii\theta})$. Classify the radii at which angular
zeros enter or leave through the cut endpoint, or form tangencies.
The elementary depth-two loss at $\rho=2$ shows that a universal
continuation without transitions is impossible. The higher-depth
problem may distinguish endpoint transitions from interior folds.

\subsection{Colored arguments and arithmetic identities}
For nontrivial inner colors, the real positive seed or the real signed
kernel can disappear. Identify color patterns admitting a real
character projection with controlled sign variation. This would be a
meaningful route toward analogous zero theorems for cyclotomic multiple
polylogarithms. Separately, pursue the manuscript's frozen $S_6$ and
$S_8$ candidates through exact functional-equation certificates. The
all-depth evaluator can supply checks, but a numerical enclosure
containing zero is not the missing arithmetic proof.

\subsection{A finite-to-analytic formalization path}
The first formal targets are the two operator moment identities,
finite polynomial cancellation, rational partial fractions with repeated
poles, and the exact evaluator's remainder arithmetic. The analytic
layer can then isolate Rolle allocation, Hausdorff-measure uniqueness,
positive-transform sectors, and the implicit-function argument.
None of the current code is represented as a Lean or other
proof-assistant formalization. A future formalization should record
exactly where classical compactness and analytic continuation enter,
rather than labeling the numerical replay as a substitute.
